% ============================================================
%  Chapitre 3 — Compréhension et préparation des données
% ============================================================
\chapter{Compréhension et préparation des données}
\label{ch:donnees}

\section{Collecte des données}

Le corpus utilisé provient du jeu de données public \emph{Company Documents
Dataset}, disponible sur la plateforme Kaggle. Il rassemble des documents
d'entreprise générés à partir de la base de démonstration \emph{Northwind}, un
jeu de données commercial classique décrivant l'activité d'un distributeur
agroalimentaire fictif : commandes, clients, produits, fournisseurs et
employés.

Ce corpus a été retenu pour trois raisons. Il est \textbf{librement
redistribuable}, ce qui autorise sa publication dans un dépôt de code. Il est
\textbf{réaliste} dans sa structure : documents hétérogènes, générés
automatiquement à partir d'un système d'information, tels qu'on en trouve
effectivement dans les archives d'une entreprise. Il est enfin d'un
\textbf{volume significatif} --- plusieurs milliers de fichiers --- permettant
de valider le passage à l'échelle du pipeline d'ingestion.

Les fichiers, tous au format PDF, ont été déposés dans le répertoire
\texttt{data/raw/} du projet, qui constitue le point d'entrée unique du pipeline
d'ingestion.

\section{Description des données}

\subsection{Composition du corpus}

Le corpus rassemble \textbf{2\,676 fichiers PDF} répartis en quatre familles
documentaires, présentées au tableau~\ref{tab:corpus}.

\begin{table}[htbp]
\centering
\caption{Composition du corpus documentaire}
\label{tab:corpus}
\small
\begin{tabularx}{\textwidth}{@{}p{3.5cm}rrX@{}}
\toprule
\textbf{Famille} & \textbf{Fichiers} & \textbf{Part} & \textbf{Contenu} \\
\midrule
Factures & 830 & 31,0~\% &
Identifiant de commande, coordonnées client, lignes de produits, montant total \\
Bons de commande & 830 & 31,0~\% &
Identifiant, date, nom du client, produits commandés \\
Commandes & 809 & 30,2~\% &
Détails d'expédition, client, employé, transporteur, lignes de produits \\
Rapports de stock & 207 & 7,8~\% &
État mensuel des stocks par catégorie et par produit \\
\midrule
\textbf{Total} & \textbf{2\,676} & \textbf{100~\%} & \\
\bottomrule
\end{tabularx}
\end{table}

\subsection{Analyse statistique}

Une analyse systématique de la longueur du texte extrait a été menée sur
l'intégralité du corpus. Ses résultats figurent au tableau~\ref{tab:stats}.

\begin{table}[htbp]
\centering
\caption{Statistiques de longueur du texte extrait (en caractères)}
\label{tab:stats}
\small
\begin{tabular}{@{}lrrrrrr@{}}
\toprule
\textbf{Famille} & \textbf{Fichiers} & \textbf{Moyenne} & \textbf{Médiane} &
\textbf{Min} & \textbf{Max} & \textbf{Pages} \\
\midrule
Factures            & 830 & 387 & 384 & 315 & 1\,029 & 1,00 \\
Bons de commande    & 830 & 210 & 207 & 152 & 820    & 1,00 \\
Commandes           & 809 & 911 & 856 & 609 & 1\,494 & 1,84 \\
Rapports de stock   & 207 & 475 & 302 & 141 & 2\,592 & 1,00 \\
\midrule
\textbf{Ensemble}   & \textbf{2\,676} & \textbf{498} & --- & 141 & 2\,592 & 1,25 \\
\bottomrule
\end{tabular}
\end{table}

Trois enseignements se dégagent de ces mesures.

\textbf{Le corpus est de volume textuel modeste.} L'ensemble représente
1\,332\,652 caractères, soit environ 1,3~Mo de texte brut pour 2\,676 fichiers.
La taille moyenne d'un document est de 498 caractères, ce qui correspond à un
paragraphe.

\textbf{L'extraction textuelle est intégralement réussie.} Aucun des 2\,676
fichiers n'a produit un texte vide. Cette observation est importante : elle
signifie que tous les PDF du corpus sont \emph{nativement textuels} et non des
images numérisées. Aucune reconnaissance optique de caractères (OCR) n'a donc
été nécessaire, ce qui aurait considérablement alourdi le pipeline. Cette
propriété favorable ne serait vraisemblablement pas vérifiée sur un corpus
d'entreprise réel, qui comporte souvent des documents scannés.

\textbf{La dispersion des longueurs est forte entre familles.} Un bon de
commande compte en moyenne 210 caractères, une commande 911, soit un rapport de
un à quatre. Les rapports de stock présentent la plus forte variance
(médiane 302, maximum 2\,592) : certains mois ne comptent que quelques
références en stock, d'autres plusieurs dizaines. Cette hétérogénéité a une
conséquence directe sur la stratégie de découpage, analysée à la
section~\ref{sec:chunking}.

\section{Exploration des données}

\subsection{Structure des documents}

L'examen du texte extrait révèle que les quatre familles suivent chacune un
gabarit régulier, hérité de leur génération automatique. Le
listing~\ref{lst:facture} présente une facture représentative.

\begin{lstlisting}[style=plain,caption={Texte extrait d'une facture (\texttt{invoice\_10248.pdf})},label={lst:facture}]
Invoice
Order ID: 10248
Customer ID: VINET
Order Date: 2016-07-04
Customer Details:
Contact Name: Paul Henriot
Address: 59 rue de l'Abbaye
City: Reims
Postal Code: 51100
Country: France
Phone: 26.47.15.10
Product Details:
Product ID Product Name Quantity Unit Price
11 Queso Cabrales 12 14.0
42 Singaporean Hokkien Fried Mee 10 9.8
72 Mozzarella di Giovanni 5 34.8
  TotalPrice 440.0
\end{lstlisting}

Un rapport de stock adopte une structure tabulaire différente, où chaque cellule
est extraite sur une ligne distincte :

\begin{lstlisting}[style=plain,caption={Extrait d'un rapport de stock (\texttt{StockReport\_2016-07.pdf})}]
Stock Report for 2016-07
Category / Product / Units Sold / Units in Stock / Unit Price
Beverages      Chang                 105   17   19
Beverages      Chartreuse verte       48   69   18
Condiments     Chef Anton's Gumbo Mix 77    0    21.35
Confections    Pavlova                95   29   17.45
\end{lstlisting}

\subsection{Nature transactionnelle du corpus et limite identifiée}
\label{sec:limite-corpus}

L'exploration met en évidence une caractéristique déterminante du corpus, qu'il
importe d'exposer clairement car elle conditionne l'interprétation des résultats
du chapitre~\ref{ch:evaluation}.

\textbf{Les 2\,676 documents sont exclusivement de nature transactionnelle.} Il
s'agit d'enregistrements structurés --- factures, bons de commande, commandes,
états de stock --- c'est-à-dire, fonctionnellement, d'une base de données
relationnelle exportée sous forme de fichiers PDF. Le corpus ne contient
\emph{aucune} procédure, \emph{aucun} manuel technique, \emph{aucune} note
d'organisation ni base de connaissances rédigée.

Cette composition a une conséquence sur le type de questions auxquelles le
système peut répondre. Il traite naturellement les interrogations factuelles
portant sur un enregistrement précis --- « quels produits contient la commande
10590 ? », « quel est le montant total de la facture 10248 ? » --- mais il ne
peut évidemment pas répondre à une question procédurale telle que « quelle est
la procédure de remboursement des notes de frais ? », faute de document la
décrivant.

\begin{encadre}
\textbf{Analyse critique.} Une architecture RAG donne son plein rendement sur
des documents \emph{rédigés}, où l'information est diffuse dans un texte suivi
et où la recherche sémantique apporte une valeur que la recherche exacte ne
peut fournir. Sur des enregistrements structurés, une requête SQL sur la base
d'origine serait plus rapide, plus exacte et moins coûteuse. Le corpus retenu
permet donc de \emph{valider techniquement} l'ensemble du pipeline --- ingestion,
vectorisation, recherche, génération, traçabilité --- mais il sous-exploite
l'intérêt propre de l'approche. L'enrichissement du corpus par des documents
rédigés constitue la première perspective d'amélioration retenue en conclusion.
\end{encadre}

\section{Préparation des données}

\subsection{Chargement et normalisation}

Le module de chargement prend en charge cinq formats --- PDF, DOCX, TXT, MD et
CSV --- et les normalise en objets \texttt{Document} de LangChain. Bien que le
corpus actuel soit exclusivement au format PDF, cette polyvalence a été
implémentée dès l'origine afin que le système accepte sans modification un
corpus d'entreprise réel, nécessairement hétérogène.

Chaque document est enrichi de quatre métadonnées, dont dépend directement la
capacité du système à citer ses sources :

\begin{lstlisting}[style=python,caption={Enrichissement en métadonnées (\texttt{src/ingestion/loader.py})}]
metadata = {
    "source":      str(path),                          # chemin complet
    "filename":    path.name,                          # nom affiche a l'utilisateur
    "file_type":   ext.lstrip("."),                    # pdf, docx, txt...
    "ingested_at": datetime.now(timezone.utc).isoformat(),
}
return Document(page_content=content, metadata=metadata)
\end{lstlisting}

Le chargement est \emph{tolérant aux pannes} : un fichier illisible ou dépourvu
de texte extractible est journalisé puis ignoré, sans interrompre le traitement
des 2\,675 autres. Sur un corpus de cette taille, l'échec global sur un seul
fichier défectueux serait inacceptable en exploitation.

\subsection{Découpage en segments}
\label{sec:chunking}

Le découpage (\eng{chunking}) est l'opération qui divise un document en segments
de taille compatible avec la fenêtre de contexte du modèle de vectorisation. Il
est réalisé par un \texttt{RecursiveCharacterTextSplitter} paramétré à
1\,000 caractères par segment et 150 caractères de chevauchement.

Ce découpeur est dit \emph{récursif} car il tente de couper aux frontières
naturelles du texte, par ordre de préférence décroissant : d'abord aux sauts de
paragraphe (\texttt{\textbackslash n\textbackslash n}), puis aux sauts de ligne,
puis aux fins de phrase, puis aux espaces, et en dernier recours à un caractère
arbitraire. L'objectif est de préserver au maximum l'unité sémantique de chaque
segment.

Le chevauchement de 150 caractères garantit qu'une information située à la
frontière de deux segments reste intégralement présente dans au moins l'un des
deux. Sans lui, une ligne de produit coupée en son milieu deviendrait
irrécupérable pour la recherche.

Chaque segment reçoit enfin un identifiant lisible reliant le segment à son
document d'origine, indispensable à la traçabilité :

\begin{lstlisting}[style=python,caption={Attribution d'un identifiant de segment (\texttt{src/ingestion/chunker.py})}]
counters: dict[str, int] = {}
for chunk in chunks:
    src = chunk.metadata.get("source", "unknown")
    counters[src] = counters.get(src, 0) + 1
    chunk.metadata["chunk_id"] = (
        f"{chunk.metadata.get('filename', src)}::chunk_{counters[src]}"
    )
\end{lstlisting}

Un segment est ainsi identifié par une chaîne du type
\texttt{invoice\_10248.pdf::chunk\_1}, directement affichable dans l'interface.

\subsection{Effet réel du découpage sur ce corpus}

L'ingestion des 2\,676 documents a produit \textbf{2\,929 segments}, soit un
ratio de 1,09 segment par document. Ce chiffre appelle une analyse.

La longueur moyenne d'un document étant de 498 caractères pour un paramètre
\texttt{chunk\_size} de 1\,000, \textbf{la grande majorité des documents tient
intégralement dans un seul segment}. Seuls 253 segments supplémentaires ont été
produits, correspondant essentiellement aux documents de la famille
« Commandes » --- dont la longueur moyenne de 911 caractères et le maximum de
1\,494 dépassent le seuil --- ainsi qu'aux rapports de stock les plus fournis.

\begin{encadre}
\textbf{Conséquence méthodologique.} Sur ce corpus, le paramètre
\texttt{chunk\_size} est \emph{quasiment inopérant} : le réglage du découpage,
qui constitue habituellement l'un des principaux leviers d'optimisation d'un
système RAG, n'a ici pratiquement aucun effet mesurable. Chaque segment
correspond en pratique à un document entier. Cette situation présente un
avantage --- le contexte transmis au modèle est toujours cohérent et complet,
jamais tronqué --- mais elle prive l'étude d'un axe d'optimisation. Sur un corpus
de documents rédigés, où un manuel de cinquante pages produirait plusieurs
centaines de segments, ce paramètre redeviendrait déterminant.
\end{encadre}

\subsection{Récapitulatif du pipeline d'ingestion}

Le tableau~\ref{tab:ingestion} résume la transformation opérée.

\begin{table}[htbp]
\centering
\caption{Bilan du pipeline d'ingestion}
\label{tab:ingestion}
\small
\begin{tabularx}{\textwidth}{@{}p{6.2cm}X@{}}
\toprule
\textbf{Étape} & \textbf{Résultat} \\
\midrule
Fichiers présents dans \texttt{data/raw/} & 2\,676 PDF \\
Fichiers chargés avec succès & 2\,676 (100~\%) \\
Fichiers sans texte extractible & 0 \\
Volume textuel total & 1\,332\,652 caractères \\
Longueur moyenne par document & 498 caractères \\
Segments produits après découpage & 2\,929 \\
Ratio segments / document & 1,09 \\
Segments indexés dans ChromaDB & 2\,929 \\
\bottomrule
\end{tabularx}
\end{table}

L'ingestion complète est déclenchée par une commande unique, ou par appel à
l'\eng{endpoint} \texttt{POST /ingest} de l'API :

\begin{lstlisting}[style=bash,caption={Lancement de l'ingestion}]
python scripts/ingest_documents.py --source data/raw --reset
\end{lstlisting}

L'option \texttt{-{}-reset} supprime la collection existante avant
ré-indexation, garantissant l'absence de doublons lors des ingestions
successives.

\section{Conclusion}

Ce chapitre a décrit le corpus de 2\,676 documents, son exploration statistique
et sa préparation. Nous avons établi que l'extraction textuelle est
intégralement réussie, que le volume total reste modeste (1,3~Mo de texte), et
que le découpage produit 2\,929 segments. Deux limites ont été identifiées et
documentées : la nature exclusivement transactionnelle du corpus, qui
sous-exploite l'intérêt propre d'une architecture RAG, et le caractère
quasi inopérant du paramètre de découpage sur des documents aussi courts. Ces
constats seront repris lors de l'interprétation des résultats. Le chapitre
suivant présente l'architecture technique construite sur cette base.
